\documentclass[10pt,respuestas,a4]{aleph-examen}
%\documentclass[10pt,a4]{aleph-examen}
% Intercambiar el comentario de las primeras lineas
% para mostrar/ocultar respuestas

% Se recomienda leer la documentación de esta
% clase en https://github.com/alephsub0/LaTeX_aleph-examen

% -- Paquetes
\usepackage{aleph-comandos}

% -- Datos del examen
\institucion{Proyecto Alephsub0}
% \carrera{}
\asignatura{Introducción al cálculo}
\tema{Examen no. 1: Lógica y Números reales}
\autor{Andrés Merino - Mario Cueva}
\fecha{Semestre 2023-1}

\logouno[4.5cm]{Logos/LogoAlephsub0-02}
\definecolor{colortext}{HTML}{0030A1}
\definecolor{colordef}{HTML}{0030A1}
\fuente{montserrat}

% -- Comandos extra
% \geometry{margin=2cm} % - Para cambiar los márgenes

%%%%%%%%%%%%%%%%%%%%%%%%%%%%%%%%%%%%%%%%
%%%%%%%%%%%%%%%%%%%%%%%%%%%%%%%%%%%%%%%%
%%%%%%%%%%%%%%%%%%%%%%%%%%%%%%%%%%%%%%%%
\begin{document}

\encabezado

\begin{preguntas}
%%%%%%%%%%%%%%%%%%%%%%%%%%%%%%%%%%%%%%%%
%%%%%%%%%%%%%%%%%%%%%%%%%%%%%%%%%%%%%%%%
%%%%%%%%%%%%%%%%%%%%%%%%%%%%%%%%%%%%%%%%
\item
    La Sonda Solar Parker es una misión de la NASA cuyo objetivo es acercarse a la corona del Sol. Suponga que la sonda detecta que en un punto determinado del espacio, la temperatura en la dirección $(1,1,1)$ disminuye a razón de dos unidades, en la dirección $(1,0,1)$ aumenta a razón de una unidad y en la dirección $(0,-1,1)$ disminuye en razón de una unidad. Si, además de esto, la sonda detecta que está en un lugar donde la temperatura puede afectar su funcionamiento y debe alejarse de tal manera que la temperatura disminuya a la mayor razón, ¿qué dirección debe seguir la sonda?\puntaje{2.0}

\begin{respuesta}
    Sean $\func{f}{\R^3}{\R}$ la función que modela la temperatura y $a\in\R^3$ el punto donde se encuentra la sonda. Se conoce que
    \[
        f'(a;(1,1,1)) = -2,
        \qquad
        f'(a;(1,0,1)) = 1
        \texty
        f'(a;(0,-1,1)) = -1.
    \]
    Se busca la dirección de máximo decrecimiento de la función, la cual es contraria al gradiente, por lo tanto, debemos calcular $\nabla f(a)$. Sabemos que
    \[
        f'(a;(1,1,1)) = \nabla f(a) \cdot (1,1,1),
        \qquad
        f'(a;(1,0,1)) = \nabla f(a) \cdot (1,0,1)
    \]
    t
    \[
        f'(a;(0,-1,1)) = \nabla f(a) \cdot (0,-1,1),
    \]
    por lo tanto
    \[
        \nabla f(a) \cdot (1,1,1) = -2,
        \qquad
        \nabla f(a) \cdot (1,0,1) = 1
        \texty
        \nabla f(a) \cdot (0,-1,1) = -1
    \]
    Si tomamos $\nabla f(a) = (u,v,w)$, tenemos que
    \[
        -2 = \nabla f(a) \cdot (1,1,1) = (u,v,w) \cdot (1,1,1) = u + v + w,
    \]
    \[
        1 = \nabla f(a) \cdot (1,0,1) = (u,v,w) \cdot (1,1,1) = u + w
    \]
    y
    \[
        -1 = \nabla f(a) \cdot (0,-1,1) = (u,v,w) \cdot (-1,1) = -v + w,
    \]
    por lo tanto
    \[
        u = 5
        \qquad
        v = -3
        \texty
        w = -4.
    \]
    Con esto, se tiene que 
    \[
        \nabla f(a) = (5,-3,-4),
    \]
    por lo tanto, la dirección que debe seguir la sonda es $(-5,3,4)$.
\end{respuesta}


%%%%%%%%%%%%%%%%%%%%%%%%%%%%%%%%%%%%%%%%
%%%%%%%%%%%%%%%%%%%%%%%%%%%%%%%%%%%%%%%%
%%%%%%%%%%%%%%%%%%%%%%%%%%%%%%%%%%%%%%%%
\item
    \begin{enumerate}
    \item 
        Utilizando el Segundo Teorema Fundamental del Cálculo calcular $\dint_0^2 x\cos(\pi x^2)\,dx$.\puntaje{1.0}
    \item
        Utilizando el Primer Teorema Fundamental del Cálculo determinar la derivada de la función definida por $F(x)= \dint_1^{x^2} \frac{\sen(t)}{t} \,dt$, para $x>0$.\puntaje{1.0}
    \end{enumerate}

\begin{respuesta}
    \begin{enumerate}[leftmargin=*]
    \item 
        Tomemos el cambio de variable
        \[
            u = \pi x^2\qquad
            du = 2\pi x\, dx
        \]
        con lo cual, tenemos
    	\[
    	    \begin{array}{c|c}
    	        x & u\\ \hline
    	        0 & 0 \\
    	        2 & 4\pi
    	    \end{array}
    	\]
    	Por lo tanto
    	\begin{align*}
    	    \int_0^2 x\cos(\pi x^2)\,dx
    	    & = \frac{1}{2\pi}\int_0^{4\pi} \cos(u)\,du \\
    	    & = \frac{1}{2\pi}\left( \sen(u) \right)\bigg|_0^{4\pi}\\
    	    & = \frac{1}{2\pi}\left(\sen(4\pi)-\sen(0)\right)=0.
    	\end{align*}
    \item
        Dado que la función definida por 
        \[
            t\mapsto \frac{\sen(t)}{t}
        \]
        es continua entre $1$ y $x^2$ y la función $x\to x^2$ es derivable, se tiene que
        \[
            F'(x) 
            = \frac{d}{dx}\left(\int_0^{x^2} \frac{\sen(t)}{t}\, dt\right)
            = \frac{\sen(x^2)}{x^2}(x^2)'
            = \frac{\sen(x^2)}{x^2}(2x)
            = \frac{2\sen(x^2)}{x}.\qedhere
        \]
    \end{enumerate}    
    
\end{respuesta}

%%%%%%%%%%%%%%%%%%%%%%%%%%%%%%%%%%%%%%%%
%%%%%%%%%%%%%%%%%%%%%%%%%%%%%%%%%%%%%%%%
%%%%%%%%%%%%%%%%%%%%%%%%%%%%%%%%%%%%%%%%
\item
    La derivada de la función definida por $f(x)=x^2$ es
    \opciones{$f'(x)=x^2$}{$f'(x)=x$}{$f'(x)=2x$}{$f'(x)=2x^2$}

\begin{respuesta}
    La opción correcta es la \textit{c)}.
\end{respuesta}
    

%%%%%%%%%%%%%%%%%%%%%%%%%%%%%%%%%%%%%%%%
%%%%%%%%%%%%%%%%%%%%%%%%%%%%%%%%%%%%%%%%
%%%%%%%%%%%%%%%%%%%%%%%%%%%%%%%%%%%%%%%%
\item
    La derivada de la función definida por $f(x)=x^2$ es
    
    \opcionesl{$f'(x)=x^2$}{$f'(x)=x$}{$f'(x)=2x$}{$f'(x)=2x^2$}

\begin{respuesta}
    La opción correcta es la \textit{c)}.
\end{respuesta}

\end{preguntas}

\end{document}